\printedpage{34}
% head: A. Grothendieck : Dérivateurs, Ch. XV, Édition des Universités de Montpellier II et Paris VII
% foot: Version 5 novembre 2010
the inclusion
\[
\underbrace{X'_{i_0}}_{\in\, \Filt^\pi} \;\longrightarrow\; X \;\simeq\; \varinjlim X_i \;.
\]
Since $X'_{i_0} \in \Filt^\pi$, hence is $\pi$-accessible, there is an $i_1 > i_0$ such that $X'_{i_0} \longrightarrow X$ factors through
\[
X'_{i_0} \;\longrightarrow\; X_{i_1} \;.
\]
Applying this same result to $X_{i_1}$, $X'_{i_1}$ \textit{etc.}, we find an infinite sequence of morphisms
\[
\begin{tikzcd}[column sep=small]
X_{i_0} \arrow[dr] & & X_{i_1} \arrow[dr] & & X_{i_2} \arrow[dr] & \\
& X'_{i_0} \arrow[ur] & & X'_{i_1} \arrow[ur] & & X'_{i_2} \quad \cdots
\end{tikzcd}
\]
such that the morphisms composed two at a time are, on the upper line, the transitions
\[
X_{i_m} \longrightarrow X_{i_n} \qquad (\; i_m \;<\; i_n \;) \;,
\]
and on the lower line, the inclusions
\[
X'_{i_m} \hookrightarrow X'_{i_n} \qquad (\; \hookrightarrow X \;) \;.
\]
Let $i_\omega = \sup_n i_n$ in $I$, hence
\[
X_{i_\omega} \;=\; \varinjlim_{n} X_{i_n} \;\simeq\; \varinjlim_{n} X'_{i_n} \;\simeq\; X'_{i_\omega} \;,
\]
so $X_{i_\omega} \longrightarrow X$ is a monomorphism.

\msp{page 57}

The last assertion of 5.7 follows at once from the fact that a filtering direct limit of monomorphisms in $\M$ is a monomorphism (condition $(e)$, p.~46).

We can make proposition 5.6 obtained in this way more precise:

\textbf{Proposition 5.7}. \texttt{[Twice 5.7.]} \textit{Let $\M$ be a category satisfying conditions} (a), (b), (c), (d) (p.~44, 45) and (e) (p.~46). \textit{Let $\M$ be a subset of} $\Ob\M$\textit{, stable under isomorphism, $\pi$ a cardinal such that the objects of} $\Filt^\pi$ \textit{are $\pi$-accessible. The following conditions are equivalent:} % [sic] Soit $\M$ (calligraphique) -> Soit $M$, la partie de Ob M est notée M partout ailleurs

\begin{enumerate}[label={}]
\item[a)] \textit{$M$ is $\pi$-accessible, and stable under filtering $\varinjlim$.}
\item[b)] \textit{$M$ is $\pi$-accessible, and $M \cap \Filt^\pi$ is stable under filtering $\varinjlim$ of cardinality $\leq \pi$.}
\item[c)] \textit{There is a subset $M^\pi$ of} $\Filt^\pi$\textit{, stable under filtering $\varinjlim$ of cardinality $\leq \pi$ and such that for every $X \in \Ob\M$, we have $X \in M$ if and only if the set $J_M$ of the $Z_\alpha \in J$ such that $Z_\alpha \in M^\pi$ is cofinal in $J$.}
\end{enumerate}

(\textit{Moreover, under condition} (c), \textit{we have}
\printedpage{35}
% head: A. Grothendieck : Dérivateurs, Ch. XV, Édition des Universités de Montpellier II et Paris VII
% foot: Version 5 novembre 2010
\msp{page 58}
$M^\pi = M \cap \Filt^\pi$.)
\begin{enumerate}[label={}]
\item[c$'$)] \textit{There is a subset $M_0$ of} $\Filt^\pi$ \textit{such that $M$ is the smallest subset of $\M$ stable under filtering direct limits.} (\textit{Moreover, we can then take $M_0$} small.)\footnote[22]{\textbf{N.B.} We have
\[
\begin{tikzcd}[ampersand replacement=\&]
(a) \arrow[r, Leftrightarrow] \arrow[d, Rightarrow] \& (b) \arrow[d, Rightarrow] \\
(c') \arrow[r, Rightarrow, "?"] \& (c)
\end{tikzcd}
\quad ,
\]
without conditions on $\pi$ (except $\pi \geq \pi_0$).}.
\end{enumerate}

\textbf{Corollary 5.8.} \textit{The maps}
\[
M \;\longmapsto\; M^\pi = M \cap \Filt^\pi \;,
\]
\[
M^\pi \quad\longmapsto\quad \parbox[t]{0.5\textwidth}{\textit{the set of $X \in \Ob\M$ such that the set $\Ss_\pi^{M^\pi}(X)$ of the $Z_\alpha \in \Ss_\pi(X)$ which are in $M^\pi$ is cofinal} (= \textit{smallest subset of $\Ob\M$ stable under filtering $\varinjlim$})}
\]
\textit{establish bijections, inverse to one another, between the set of $M \subset \Ob\M$ satisfying the equivalent conditions of proposition} 5.7\textit{, and the set of subsets $M^\pi$ of} $\Filt^\pi$ \textit{stable under filtering $\varinjlim$ of cardinality $\leq \pi$.}

\textsc{Proof of} 5.7. The implication $(b) \Longrightarrow (c)$ is tautological, in view of the definition of the $\pi$-accessibility of $M$, and of proposition 5.3. And $(a) \Longrightarrow (b)$ comes from the stability of $\Filt^\pi$ under filtering $\varinjlim$ of cardinality $\leq \pi$ (condition $(d)$).

It therefore remains to prove $(c) \Longrightarrow (a)$. (It is here that we must use the special hypothesis on $\pi$.) Once this is seen, we shall know the equivalence of $(a)$, $(b)$, $(c)$, and on the other hand, it is immediate that $(a) \Longrightarrow (c')$, by
\msp{page 59}
taking $M_0 = M \cap \Filt^\pi$. On the other hand, $(c') \Longrightarrow (c)$ (without hypotheses on $\pi\,$?), since we can suppose that the subset $M_0$ of $\Filt^\pi$ is stable under filtering $\varinjlim$ of cardinality $\leq \pi$. But then, knowing the equivalence of $(a)$, $(b)$ $(c)$, whence the corollary at once, it follows that the envelope stable under filtering $\varinjlim$ of $M_0$ in $\M$ is none other than the accessible subset deduced from $M_0$ that is in question in $(c)$, and we win. % [sic] $(b)$ $(c)$ -> $(b)$, $(c)$, virgule manquante

So it all comes down to proving that, if all the objects of $\Filt^\pi$ are accessible, then $(c)$ implies $(a)$. Let us first show that $(c)$ implies the membership criterion (5.3) (p.~51), so by proposition 5.6 ``revised and corrected'' in note 21, it already follows that $(c) \Longrightarrow (b)$, and it will remain to establish $(b) \Longrightarrow (a)$, \textit{i.e.} that $(b)$ implies the stability of $M$ under arbitrary filtering $\varinjlim$ (not only large relative to $\pi$).

That $(c)$ implies (5.3) will follow from the following lemma, applied to $I = \Ss_\pi(X)$:

\textbf{Lemma 5.8.} \texttt{[Twice 5.8.]} \textit{Let $I$ be an ordered set that is} $\pi$-adapted, i.e. \textit{filtering and such that every filtering subset $K$ of $I$, of cardinality $\leq \pi$, has a supremum $i_K$ in $I$. I say that a subset $J$ of $I$ is} $\pi$-cofinal, \textit{if it is cofinal, and if for every filtering subset $K$ of $J$ of cardinality $\leq \pi$, its supremum $i_K$ in $I$ is}
\printedpage{36}
% head: A. Grothendieck : Dérivateurs, Ch. XV, Édition des Universités de Montpellier II et Paris VII
% foot: Version 5 novembre 2010
\msp{page 60}
\textit{in $J$.} (\textit{It follows, then, that $J$ too is $\pi$-adapted, for the induced order structure.}) \textit{This said, the intersection of two $\pi$-cofinal subsets of $I$ is again $\pi$-cofinal.}

It is enough to see that it is cofinal, for its stability under suprema in $I$ of filtering subsets of cardinality $\leq \pi$ is tautological (if it is true for some $J_\alpha$, it is true for $\bigcap_\alpha J_\alpha$). Let $i_0 \in I$, and let us construct an element of $J \cap J'$ that dominates $i_0$. We construct step by step a double sequence
\[
\begin{tikzcd}
i_1 \arrow[dr] & & i_2 \arrow[dr] & & \cdots \\
& i'_1 \arrow[ur] & & i'_2 \arrow[ur] &
\end{tikzcd}
\]
with the $i_\alpha$ in $J$, the $i'_\alpha$ in $J'$, and $i_1 \geq i_0$. We thus have
\[
\sup i_\alpha \;=\; \sup i'_\alpha \;,
\]
and if $i$ is the common value of the two sides, we have $i \in J \cap J'$.

This seen, to prove the validity of (5.3), we note that the subset $J' = \Ss_\pi^{(M^\pi)}(X)$ of $I = \Ss_\pi(X)$, consisting of the $Z_\alpha \subset X$ such that $Z_\alpha \in M^\pi$, is stable under suprema of filtering subsets of cardinality $\leq \pi$, since $M^\pi$ is stable
\msp{page 61}
under $\varinjlim$ of cardinality $\leq \pi$. So if $J$ is a $\pi$-cofinal subset of $\Ss_\pi(X)$, $J \cap J'$ is cofinal in $J$, \textit{i.e.} in $I$, if and only if $J'$ is cofinal in $I$ (\textit{i.e.} if and only if $X \in M$, by condition $(c)$). For $J \cap J'$ cofinal $\Longrightarrow$ $J'$ cofinal tautologically, and $J'$ cofinal implies $J'$ $\pi$-cofinal, hence $J \cap J'$ $\pi$-cofinal by lemma 5.8, hence cofinal.

It remains only to prove that $(b) \Longrightarrow (a)$, \textit{i.e.} the stability of $M$ under arbitrary filtering direct limits, assuming $(b)$. Here again, this should not use the special hypothesis on $\pi$, nor even the exactness hypothesis $(e)$ on $\M$, but be a general story of accessibility:

\textbf{Proposition 5.9.} \textit{Let $\M$ be a category satisfying} (a), (b), (c), (d$_0$) (pages 44, 45), \textit{$\pi \geq \pi_0$ a cardinal, $M$ a $\pi$-accessible subset of} $\Ob\M$, i.e. \textit{satisfying condition $b_\pi$ of} XIII, p.~88. \textit{Then $M$ is stable under filtering $\varinjlim$ if and only if $M^\pi \overset{\text{def}}{=} M \cap \Filt^\pi$ is stable under filtering $\varinjlim$ of cardinality $\leq \pi$.}

I don't feel like spelling out the proof here, having the painful feeling of forever redoing the same proof (here and in XIII). I shall have to come back to this in a more settled situation.

\msp{page 62}

I am going to sketch the proof all the same. Let $X = \varinjlim\limits_{I} X_i$, $I$ filtering, the $X_i$ in $M$, we must prove that $X \in M$ (assuming the stability hypothesis on $M^\pi = M \cap \Filt^\pi$).
\printedpage{37}
% head: A. Grothendieck : Dérivateurs, Ch. XV, Édition des Universités de Montpellier II et Paris VII
% foot: Version 5 novembre 2010
\[
\begin{tikzcd}
& \Ci \arrow[r] \arrow[d] \arrow[dl, dotted, "\mathfrak{z}"'] \arrow[dr, phantom, "\text{\footnotesize cart.}"] & I \arrow[d, "\varphi"] \\
\underline{M}^\pi & \underline{M}^\pi/\underline{M} \arrow[l, dotted, "\text{\footnotesize ``source''}"] \arrow[r, "\text{\footnotesize ``target''}"'] & \underline{M}
\end{tikzcd}
\] % resolved: the PDF glyph is fraktur z; see corrections.csv, leaf 4.

$\Ci$ is a cofibered category over $I$, whose fibers are
\[
\Ci_i \;\simeq M^\pi/X_i \;,
\]
now since $X_i \in M$, we see, because of the $\pi$-accessibility condition, that $M^\pi/X_i$ is \textit{filtering.} We are going to single this out as

\textbf{Lemma 5.10.} \textit{Let $M$ be a $\pi$-accessible subset of $\M$. Then for every $Y \in M$, $M^\pi/Y$ is a category that is} filtering, \textit{and $\Ss_\pi(Y)$ is cofinal in it.}

\msp{page 63}

Let us consider the canonical functor
\[
I \;=\; \Ss_\pi(Y) \xrightarrow{\;\varphi\;} \underline{M}^\pi/Y \;=\; I' \;,
\]
where the source category $I$ is filtering. By SGA 4 I 8.1.3, this functor is cofinal if and only if it satisfies two conditions F\,1, F\,2, and then $I'$ is also filtering. Condition F\,1 is the usual condition that every $i' \in \Ob I'$ is ``dominated'' by a $\varphi(i)$, \textit{i.e.} every $Z \longrightarrow Y$, with $Z \in M^\pi$, is ``dominated'' by a \textit{monomorphism} $Z_\alpha \hookrightarrow Y$, with $Z_\alpha \in M^\pi$, which follows from the fact that it is dominated by $Z_\alpha \hookrightarrow Y$ with $Z_\alpha \in \Filt^\pi$, \textit{i.e.} $Z_\alpha \in \Ss_\pi(Y)$, and that among these $Z_\alpha$ those which are in $M^\pi$ are cofinal (since $Y \in M$). The other condition\footnote[23]{Since $I$ is filtering.} is that a double arrow $i' \overset{f'}{\underset{g'}{\rightrightarrows}} \varphi(i)$ can
\msp{page 64}
be equalized by an arrow $\varphi(h)$, with $h : i \longrightarrow j$. But in fact, for an $i' \in \Ob I'$, \textit{i.e.} a given $Z \longrightarrow Y$ ($Z \in M^\pi$), and an $i \in \Ob I$, \textit{i.e.} a given $Z_\alpha \hookrightarrow Y$ ($Z_\alpha \in M^\pi$), there is at most \textit{one} arrow in $I' = \underline{M}^\pi/Y$ from $Z \longrightarrow Y$ to $Z_\alpha \hookrightarrow Y$. So condition F\,2 is trivially satisfied (we take $h = \mathrm{id}_i$).

\textbf{Corollary 5.11.} \textit{If $(X_i)_{i \in I}$ is a filtering direct system, with the $X_i$ in $M^\pi$, then the category $\Ci = \underline{M}^\pi/\underline{M} \times_{\underline{M}} (I,\varphi)$, constructed on page} 62, \textit{is filtering, and the canonical functor} (\textit{fibrant}) \textit{$\Ci \longrightarrow I$ is cofinal.}

That $\Ci$ is filtering and $\Ci \longrightarrow I$ cofinal follows formally from the fact that $I$ is, that the fibers of $\Ci$ are, and that $\Ci$ is cofibered over $I$. (Deserves a lemma, for future reference.)

\textbf{N.B.} We can replace $\Ci$ by $\Ci_0$ whose fibers are the
\[
\Ss_\pi^{M^\pi}(X_i) \;\subset\; \underline{M}^\pi \;; \qquad \text{cofinal in } \Ci_i = \underline{M}^\pi/X_i \;,
\]
\msp{page 65}
we find a subcategory of $\Ci$ cofinal in $\Ci$, and which this time is ordered just like $I$.

\printedpage{38}
% head: A. Grothendieck : Dérivateurs, Ch. XV, Édition des Universités de Montpellier II et Paris VII
% foot: Version 5 novembre 2010
We have
\[
X_i \;=\; \varinjlim_{Z_\alpha \in \Ci_i} Z_\alpha \;,
\]
hence
\[
X \;=\; \varinjlim_{\alpha \in \Ci} Z_\alpha \;,
\]
altogether we have arrived at representing $X$ as a filtering $\varinjlim$ of elements $Z_\alpha$ which are (not only in $M$, but even) in $M^\pi$.

Let us change the notation, taking
\[
X \;=\; \varinjlim_{i \in I} X_i \;, \qquad I \text{ filtering ordered set, } X_i \in M^\pi \;\; \forall\, i \in I \;,
\]
let $I'$ be the set of \textit{filtering} subsets $J \subset I$ such that $\mathrm{card}\, J \leq \pi$, $I'$ ordered by inclusion. Then $I'$ is large relative to $\pi$, and we have a direct system
\[
(X_J)_{J \in I'} \;, \qquad X_J \;=\; \varinjlim_{i \in J} X_i \;.
\]
By the stability hypothesis on $M^\pi$ (which has not yet been used), the $X_J$ are in $M^\pi$. We therefore have
\[
X \;=\; \varinjlim_{J \in I'} X_J \;, \qquad \text{the } X_J \text{ in } M^\pi\text{, } I' \text{ large relative to } \pi \;,
\]
so $X \in M$ since $M$ is $\pi$-admissible, \quad \textit{q.e.d.}

\msp{page 66}

\textbf{Definition 5.12.}
\begin{enumerate}[label={}]
\item[1)] Let $\M$ be a category satisfying conditions $(a)$, $(b)$, $(c)$, $(d)$ (p.~44-46). Let $\pi \geq \pi_0$ be a cardinal. A subset $M$ of $\Ob\M$ is said to be \textit{L-$\pi$-accessible,} if it is $\pi$-accessible and stable under small filtering $\varinjlim$\footnote[24]{We then say that $M$ is \textit{L-stable.}}, or, what amounts to the same thing (5.9), if it is $\pi$-accessible and if $M^\pi = M \cap \Filt^\pi$ is stable under filtering $\varinjlim$ of cardinality $\leq \pi$. (\textbf{N.B.} If the objects of $\Filt^\pi$ are $\pi$-accessible, 5.8 tells us that this also means that it suffices for this that there is a subset $M_0$ of $\Filt^\pi$ such that $M$ is the L-envelope (envelope stable under filtering $\varinjlim$) of $M_0$ -- a condition which is in any case always necessary, by taking $M_0 = M^\pi$.)
\item[2)] Let $M$ be a subset of $\Ob\M$. We say that it is \textit{L-accessible,} if it satisfies one of the following four equivalent conditions:
\begin{enumerate}[label={}]
\item[\textit{a})] There is a cardinal $\pi$ such that $M$ is L-$\pi$-accessible.
\item[\textit{a}$'$)] $M$ is accessible and L-stable.
\item[\textit{b})] There is a small subset $M_0$ of $\Ob\M$, such that $M$ is the L-envelope of $M_0$\footnote[25]{This condition does not depend on the chosen cardinal filtration of $\M$.}.
\item[\textit{b}$'$)] There is a subset $M_0$ of $\M$ contained in a subset $\Filt^\pi$ (or, what amounts to the same thing, such that the full subcategory $\underline{M}_0$ of $\M$ is essentially small, \textit{i.e.} equivalent to a small category, such that $M$ is the \textit{L}-envelope of $M_0$).
\end{enumerate}
\end{enumerate}